% ch4_c §4.6–4.7 (书页 139–145 / p153–p159)
%--B-TAIL--
例如，为了计算 $\zeta$ 随温度的变化，我们利用式 (4.9) 和式 (4.10)：
\begin{align*}
\int_0^{\mathcal{E}_F} \mathcal{N}(\mathcal{E})\,d\mathcal{E} &= \int_0^{\infty} \mathcal{N}(\mathcal{E})f^0(\mathcal{E})\,d\mathcal{E} \\
&= \int_0^{\zeta} \mathcal{N}(\mathcal{E})\,d\mathcal{E} + \frac{\pi^2}{6}(kT)^2\left[\frac{\partial \mathcal{N}(\mathcal{E})}{\partial \mathcal{E}}\right]_{\mathcal{E}=\zeta} + \ldots,
\tag{4.22}
\end{align*}
它具有近似解（可通过关于 $\zeta$ 求微商来检验）
\begin{equation*}
\zeta \approx \mathcal{E}_F - \frac{\pi^2}{6}(kT)^2\left[\frac{\partial}{\partial \mathcal{E}}\ln \mathcal{N}(\mathcal{E})\right]_{\mathcal{E}=\mathcal{E}_F}.
\tag{4.23}
\end{equation*}

%FIG{82}: {图 82　高温下，费米–狄拉克 函数变为经典 Boltzmann 分布。}

我们注意到 $\zeta<\mathcal{E}_F$，只有 $T=0$ 时例外——但当 $kT\ll\mathcal{E}_F$ 时，例如常温下的金属，这一修正很小。随着温度升高，$\zeta$ 的值不得不下降，因为 费米–狄拉克 函数在电子气的简并温度（degeneracy temperature）以上——即当 $kT>\mathcal{E}_F$ 时——必须趋于经典分布
\begin{equation*}
f^0(\mathcal{E})\propto e^{-\mathcal{E}/kT}
\tag{4.24}
\end{equation*}
从式 (4.8) 的形式显然可见，这只对 $\zeta$ 以上的能量才成立；若要它对带中所有能量都成立，$\zeta$ 的值就必须位于 $\mathcal{E}=0$ 之下。

\section{半导体中载流子的统计}

在半导体中，态密度函数有一个很大的间隙。实际上，能量的积分范围被分成两段——一段到价带（valence band）顶 $\mathcal{E}_v$ 为止，另一段从导带（conduction band）底 $\mathcal{E}_c$ 起继续向上。于是，与式 (4.9) 和式 (4.10) 相当的条件是
\begin{equation*}
\int_0^{\mathcal{E}_v} \mathcal{N}(\mathcal{E})\,d\mathcal{E} = \int_0^{\mathcal{E}_v} f^0(\mathcal{E})\mathcal{N}_v(\mathcal{E})\,d\mathcal{E} + \int_{\mathcal{E}_c}^{\infty} f^0(\mathcal{E})\mathcal{N}_c(\mathcal{E})\,d\mathcal{E},
\tag{4.25}
\end{equation*}
其中 $\mathcal{N}_c(\mathcal{E})$ 与 $\mathcal{N}_v(\mathcal{E})$ 分别是导带与价带中的态密度。

我们可以把它写成
\begin{equation*}
\int_0^{\mathcal{E}_v} \{1-f^0(\mathcal{E})\}\mathcal{N}_v(\mathcal{E})\,d\mathcal{E} = \int_{\mathcal{E}_c}^{\infty} f^0(\mathcal{E})\mathcal{N}_c(\mathcal{E})\,d\mathcal{E},
\tag{4.26}
\end{equation*}
这不过是说，被激发到导带中的电子数等于价带中留下的空穴（hole）数：
\begin{equation*}
n_h = n_e.
\tag{4.27}
\end{equation*}

式 (4.26) 中的统计函数可以改写成一种有趣的形式。让我们注意
\begin{align*}
1-f^0(\mathcal{E}-\zeta) &= 1 - \frac{1}{\exp\{(\mathcal{E}-\zeta)/kT\}+1} \\
&= \frac{1}{\exp\{-(\mathcal{E}-\zeta)/kT\}+1} \\
&= f^0(\mathcal{E}_h+\zeta_h),
\tag{4.28}
\end{align*}
其中
\begin{equation*}
\mathcal{E}_h+\zeta_h = -(\mathcal{E}-\zeta).
\tag{4.29}
\end{equation*}

换言之，在电子的费米能级（Fermi level）之下距离 $|\mathcal{E}-\zeta|$ 处的能级 $\mathcal{E}$ 上找到一个空穴的概率，与在费米能级之上距离 $|\mathcal{E}-\zeta|$ 处找到一个电子的概率相同。空穴与电子一样遵从 费米–狄拉克 统计——只不过仿佛它们的能量是\emph{向下}量的。为方便起见，我们把所有电子能量都从 $\mathcal{E}_c$ 量起；用变量 $\mathcal{E}_e$ 表示 $\mathcal{E}-\mathcal{E}_c$。空穴能量则从价带顶量起：$\mathcal{E}_h=-(\mathcal{E}-\mathcal{E}_v)$。电子的费米能级位于某一点 $\mathcal{E}_e=-\zeta_e$ 处；空穴的费米能级则相应于 $\mathcal{E}_h=-\zeta_h$。于是我们的公式 (4.26) 便读作
\begin{equation*}
\int_0^{\infty} f^0(\mathcal{E}_h+\zeta_h)\mathcal{N}_v(\mathcal{E}_h)\,d\mathcal{E}_h = \int_0^{\infty} f^0(\mathcal{E}_e+\zeta_e)\mathcal{N}_c(\mathcal{E}_e)\,d\mathcal{E}_e.
\tag{4.30}
\end{equation*}

这两个函数还须满足的进一步条件是：$\zeta_e$ 与 $\zeta_h$ 事实上应当对应于同一个总的费米势（over-all Fermi potential）$\zeta$。既然它们是从彼此相距 $\mathcal{E}_{\mathrm{gap}}$ 的能量原点分别量起的，我们便有\footnote{这里看来最方便的做法，是把 $\zeta_e$ 和 $\zeta_h$ 定义成通常均为正的量。}
\begin{equation*}
\zeta_e+\zeta_h = \mathcal{E}_{\mathrm{gap}}.
\tag{4.31}
\end{equation*}

正如我们马上就要证明的，$\zeta$ 倾向于落在间隙之中，因此 $\zeta_e$ 与 $\zeta_h$ 都是正的，而且比 $kT$ 大得多。于是，在 $\mathcal{E}_e$ 与 $\mathcal{E}_h$ 为正的范围内，电子和空穴的 Fermi 函数实际上已无法与经典分布区分：
\begin{equation*}
f^0(\mathcal{E}_e+\zeta_e) \approx \exp\{-(\mathcal{E}_e+\zeta_e)/kT\}.
\tag{4.32}
\end{equation*}

%FIG{83}: {图 83　(a) 半导体的绝对能量图式。(b) “电子”与“空穴”能量图式。}

这在一定程度上简化了式 (4.30) 中积分的计算。按照 van Hove 定理（§2.5），导带底附近态密度的形式必定是
\begin{equation*}
\mathcal{N}_c(\mathcal{E}_e) = \frac{1}{2\pi^2}\left(\frac{2m_e}{\hbar^2}\right)^{3/2}\mathcal{E}_e^{1/2}.
\tag{4.33}
\end{equation*}
这实际上就是\emph{有效质量}（effective mass）为 $m_e$ 的自由电子气的结果，但一般说来这个参量是
\begin{equation*}
m_e = (m_1 m_2 m_3)^{1/2},
\tag{4.34}
\end{equation*}
其中 $m_1$、$m_2$、$m_3$ 是在 $\mathbf{k}$ 空间中围绕 $\mathcal{E}(\mathbf{k})$ 的极小值、以局部主轴为参照展开能量时的展开系数：
\begin{equation*}
\mathcal{E}_e = \frac{\hbar^2 k_1^2}{2m_1} + \frac{\hbar^2 k_2^2}{2m_2} + \frac{\hbar^2 k_3^2}{2m_3}.
\tag{4.35}
\end{equation*}
其证明与式 (2.70)–(2.72) 完全相同。

把式 (4.32) 与式 (4.33) 代入 (4.30)，我们得到
\begin{align*}
n_e &= \int_0^{\infty} \exp\{-(\mathcal{E}_e+\zeta_e)/kT\}\,\frac{1}{2\pi^2}\left(\frac{2m_e}{\hbar^2}\right)^{3/2}\mathcal{E}_e^{1/2}\,d\mathcal{E}_e \\
&= \frac{1}{2\pi^2}\left(\frac{2m_e}{\hbar^2}\right)^{3/2} e^{-\zeta_e/kT}(kT)^{3/2}\int_0^{\infty} e^{-z}z^{1/2}\,dz \\
&= 2\left(\frac{m_e kT}{2\pi\hbar^2}\right)^{3/2} e^{-\zeta_e/kT}.
\tag{4.36}
\end{align*}

在简单的模型中——例如忽略自旋–轨道分裂（§3.11）时——价带顶也将具有与式 (4.33) 相同函数形式的态密度，只是改由一个具有质量量纲的类似参量 $m_h$ 刻画。对于“在温度 $T$ 下激发的空穴”的总数，将有类似的公式：
\begin{equation*}
n_h = 2\left(\frac{m_h kT}{2\pi\hbar^2}\right)^{3/2} e^{-\zeta_h/kT},
\tag{4.37}
\end{equation*}

令 $n_e$ 与 $n_h$ 相等，并利用 (4.31)，就可以把二者分别求出。例如，将 (4.36) 与 (4.37) 相乘，得
\begin{equation*}
n_e n_h = 4\left(\frac{kT}{2\pi\hbar^2}\right)^3 (m_e m_h)^{3/2} e^{-\mathcal{E}_{\mathrm{gap}}/kT},
\tag{4.38}
\end{equation*}
此式与费米能级的位置无关，因此只要费米能级保持在带隙之内，它就相当普遍地成立。对此式取平方根，得
\begin{equation*}
n_e = n_h = 2\left(\frac{kT}{2\pi\hbar^2}\right)^{3/2}(m_e m_h)^{3/4} e^{-\mathcal{E}_{\mathrm{gap}}/2kT},
\tag{4.39}
\end{equation*}
这表明载流子（carrier）密度确实按 Boltzmann 因子变化，就好像需要能量 $\frac{1}{2}\mathcal{E}_{\mathrm{gap}}$ 才能把一个电子激活到导带中去。

由 (4.36)、(4.37) 和 (4.39)，容易分别算出 $\zeta_e$ 与 $\zeta_h$；由对称性显然可知，若 $m_e=m_h$，$\zeta$ 将落在间隙正中。只要 $\mathcal{E}_g\gg kT$，这些公式就对载流子全部由热激发产生的\emph{本征半导体}（intrinsic semiconductor）成立。

在本书中，我们把讨论主要限于纯净、有序的固体；但应当提到，大多数半导体并不很纯——实际上，我们几乎总是在掺入杂质、经过\emph{掺杂}（doped）的状态下使用它们。标准的情形是\emph{施主}（donor）杂质——一个以替位方式掺入 Ge 晶体的 As 原子。由 §4.2 的论证显然可知，只要让 As 的第 5 个电子被释放到晶体的导带之中，这种替换就可以在不损害晶格键结构的情况下实现。这样我们便制得一个 \emph{n 型}（n-type）样品；我们增加了几个电子，却没有在价带中产生任何空穴。

若每单位体积中有 $N_d$ 个这样的\emph{施主原子}，则必有
\begin{equation*}
n_e - n_h = N_d
\tag{4.40}
\end{equation*}
以此作为条件，从中求出 $n_e$、$n_h$、$\zeta_e$、$\zeta_h$ 随 $T$ 变化的函数关系。特别是，若通过添加大量施主使材料重度掺杂，电子将占压倒性的优势。这时费米能级必定落在这些\emph{多数载流子}（majority carriers）所占据的能量范围内——即靠近导带底。掺杂足够多时，这种载流子气可以达到金属的密度，从而即使在常温下分布也变成简并的。

%FIG{84}: {图 84　(a) Ge 中的施主。(b) Ge 中的受主。}

这类杂质还有一种对偶的情形。如果我们用 In——每个原子只有 3 个电子——来代替 Ge，键就不能完全填满；价带中留下一个空穴。晶体中密度为 $N_a$ 的\emph{受主}（acceptor）杂质造成一个 \emph{p 型}（p-type）样品，其中
\begin{equation*}
n_e - n_h = -N_a.
\tag{4.41}
\end{equation*}

一般地，施主与受主可能同时存在，于是在\emph{补偿}（compensated）样品中有
\begin{equation*}
n_e - n_h = N_d - N_a.
\tag{4.42}
\end{equation*}
不过，这些杂质也会把局域的杂质能级引入能隙之中（见 §6.4）；因此，当温度低到这些能级可能被相当大比例的过剩载流子占据时，统计理论就变得更复杂了。

\section{电子比热}

金属中的电子必然对比热有所贡献。为了计算这一贡献，让我们利用式 (4.21) 计算电子分布的平均能量，因为我们假定系统是高度简并的：
\begin{align*}
\overline{\mathcal{E}} &= \int \mathcal{E}f^0(\mathcal{E})\mathcal{N}(\mathcal{E})\,d\mathcal{E} \\
&= \int_0^{\zeta} \mathcal{E}\mathcal{N}(\mathcal{E})\,d\mathcal{E} + \frac{\pi^2}{6}(kT)^2\left[\frac{\partial}{\partial \mathcal{E}}\{\mathcal{E}\mathcal{N}(\mathcal{E})\}\right]_{\mathcal{E}=\zeta} + \ldots.
\tag{4.43}
\end{align*}

我们对温度求微商，同时注意到（参见式 (4.22)）$\zeta$ 是 $T$ 的函数：
\begin{align*}
C_{\mathrm{el.}} &= \frac{\partial \overline{\mathcal{E}}}{\partial T} = \zeta\mathcal{N}(\zeta)\frac{d\zeta}{dT} + \frac{\pi^2}{3}k^2T\left[\mathcal{N}(\mathcal{E}) + \zeta\frac{\partial \mathcal{N}(\mathcal{E})}{\partial \mathcal{E}}\right]_{\mathcal{E}=\zeta} + O(T^2) \\
&= \frac{\pi^2}{3}k^2T\mathcal{N}(\mathcal{E}_F) + \zeta\mathcal{N}(\zeta)\left[\frac{d\zeta}{dT} + \frac{\pi^2}{3}\frac{k^2T}{\mathcal{N}(\mathcal{E})}\frac{\partial \mathcal{N}(\mathcal{E})}{\partial \mathcal{E}}\right]_{\mathcal{E}=\zeta} \\
&= nk\,\frac{\pi^2}{3}kT\,\frac{\mathcal{N}(\mathcal{E}_F)}{n}
\tag{4.44}
\end{align*}
其中最后一步由式 (4.22) 得出。$T^2$ 等更高阶的项都可以忽略。

这是一个十分重要的结果。把式 (4.44) 与经典粒子气体的比热——譬如说 $\frac{3}{2}nk$——作一比较，量子结果要小得多。对自由电子，态密度式 (4.3) 为 $\frac{3}{2}n/\mathcal{E}_F$，于是
\begin{equation*}
C_{\mathrm{el.}}/C_{\mathrm{classical}} = \frac{\pi^2}{3}\frac{kT}{\mathcal{E}_F}.
\tag{4.45}
\end{equation*}
在普通金属中、普通温度下，这个因子约为 $1/100$。它说明了为什么金属的总比热可以由 §2.4 的晶格项相当准确地给出，也说明了为什么 Dulong–Petit 定律在高温下成立。

我们还注意到 $C_{\mathrm{el.}}$ 与 $T$ 成线性。在极低温度下，这一线性项——通常写作
\begin{equation*}
C_{\mathrm{el.}} = \gamma T,
\tag{4.46}
\end{equation*}
——可以与晶格项区分开来，后者按 $T^3$ 更快地趋于零。测量 $\gamma$ 可以直接给出关于 $\mathcal{N}(\mathcal{E}_F)$——费米能级处的态密度——的信息。例如，在过渡金属中观察到很大的 $\gamma$ 值，正如 §4.4 中所讨论的。

我们可以这样理解这条线性定律。看看 费米分布（图 85）；温度的作用是使少数电子被激发到更高的能级。但这一效应只在 $\mathcal{E}_F$ 附近 $kT$ 量级的能量范围内才是明显的。我们可以说，大约有 $kT\,\mathcal{N}(\mathcal{E}_F)$ 个电子各获得约 $kT$ 的能量。于是，整个固体获得的能量为
\begin{equation*}
\delta\overline{\mathcal{E}} \sim k^2T^2\mathcal{N}(\mathcal{E}_F).
\tag{4.47}
\end{equation*}
与此相应，
\begin{equation*}
C_{\mathrm{el.}} = \frac{\partial \delta\overline{\mathcal{E}}}{\partial T} \sim 2k^2T\mathcal{N}(\mathcal{E}_F).
\tag{4.48}
\end{equation*}
除一个数值因子外，这基本上就是我们的结果 (4.44)。

%FIG{85}: {图 85　金属中电子的热激发。}

实际上，我们要说的是：处在分布深处的电子几乎觉察不到温度的存在。它们的处境由 Pauli 原理决定；这一原理坚持要它们填满所有的能级，却不允许它们侵占彼此的状态。这并不比下面这件事更令人奇怪：离子实束缚态深处的电子不必算作对固体比热的一种单独贡献——至少在温度还没有高到足以把这些电子热激发起来之前是如此。
